\documentclass[11pt]{article}
\usepackage[T1]{fontenc}
\usepackage{lmodern}
\usepackage[margin=1in]{geometry}
\usepackage{amsmath,amssymb,amsthm,booktabs,microtype}
\usepackage[colorlinks=true,linkcolor=blue,citecolor=blue,urlcolor=blue]{hyperref}
\hypersetup{pdftitle={A conditional proper-time volume response for prepared Einstein-scalar shell data},pdfauthor={Ricardo Maldonado},pdfsubject={Working manuscript for author and specialist review}}
\pdfinfoomitdate=1
\pdftrailerid{}
\pdfsuppressptexinfo=15
\newcommand{\eps}{\varepsilon}
\newcommand{\dd}{\mathrm d}
\newcommand{\Lie}{\mathcal L}
\newcommand{\av}[1]{\left\langle #1\right\rangle}
\newcommand{\norm}[1]{\left\lVert #1\right\rVert}
\newcommand{\bQ}{\mathcal Q}
\newcommand{\vol}{\operatorname{Vol}}
\newtheorem{proposition}{Proposition}
\title{A conditional proper-time volume response for prepared Einstein--scalar shell data}
\author{Ricardo Maldonado\\\small Independent Researcher}
\date{1 November 2026}

\begin{document}
\maketitle
\begin{center}\small Working manuscript for author and specialist review\end{center}
\begin{abstract}
We calculate a conditional second variation of anchored proper-time volume
for a compactly prepared perturbation of a five-dimensional Einstein--scalar
spacetime with a pure-tension timelike shell. An explicit compact monopole
particular solution fixes second initial data that were not numerically
specified by an earlier weighted-only completion. A pure second-order gauge
transformation removes a large spatial Lie-derivative term while retaining
its scalar contribution. The complete five-component mean system is bounded
by a fresh causal contraction estimate. Clock and volume-measure corrections
give the proper scalar mean, and an exact intrinsic Gauss identity with a
positive Green kernel gives the volume and expansion response. At the fixed
proper time $\tau_*=(5/4)\rho_b10^{-6}$, outward enclosures are
$V_2/V_0\in[-5.4520,-2.6521]\times10^{16}$ and
$\Theta_2\in[-9.0208,-4.4458]\times10^{21}$ in the stated model and parameter
normalization. The conclusions assume the retained background, localized
constraint inverse, smooth shell development, and source-identification
results. They imply smaller volume and smaller positive expansion than the
background for sufficiently small nonzero parameters in each fixed family;
they give no effective parameter radius or physical radiation mechanism.
\end{abstract}

\section{Problem, model, and conditional statement}

Second-order observables depend on the nonlinear completion of first-order
Cauchy data. A fixed first packet and zero initial tensor-transverse-traceless
data do not determine the scalar second response. In particular, a weighted
constraint construction whose radial weights have not been instantiated
does not identify a unique numerical mean history. We address this problem
by prescribing a compact mean particular solution and applying the retained
constraint construction to the remaining source. This is a new preparation,
not a reconstruction of unspecified old weights.

The bulk and one-copy outward shell equations are
\begin{align}
 \operatorname{Ric}_{AB}-\phi_A\phi_B-\tfrac23\mathcal V(\phi)g_{AB}&=0,
 &\Box_g\phi-\mathcal V'(\phi)&=0, \label{eq:model}\\
 K^{\rm out}_{ab}&=\tfrac16\sigma(\phi)\gamma_{ab},
 &n_{\rm out}\phi&=-\tfrac12\sigma'(\phi). \label{eq:shell}
\end{align}
Two identical interior copies use their respective outward normals. The
retained model is
\begin{gather}
 W(\phi)=1-\phi+\phi^3/3,\qquad
 \mathcal V=\tfrac12(W')^2-\tfrac23W^2,\qquad
 \sigma=2W+\delta(1+c_\sigma\phi), \nonumber\\
 \delta=10^{-3},\qquad c_\sigma=96714321425692/161839258930731.
 \label{eq:potential}
\end{gather}
The selected near-$\phi=1$ regular background, with its retained shell
data, is
\begin{equation}
 g_0=\dd y^2+\rho(y)^2[-\dd t^2+\cosh^2t\,\kappa],\qquad
 a=\rho_y/\rho,\quad p=\phi_{0,y}, \label{eq:background}
\end{equation}
where $\kappa$ is the unit round metric on $S^3$. Put
$\ell=10^{-6}$, $x=t/\ell$, $r=z/\ell$, and $\dd y=-\rho\,\dd z$.
Thus $\partial_r=-\ell\rho\partial_y$. The shell radius is $\rho_b$.
The coordinate $t$ is not proper time; on the background shell
$\tau=\rho_b t$.

All subscript-$2$ quantities are second parameter derivatives:
\begin{equation}
 g_\eps=g_0+\eps q+\tfrac12\eps^2 k+o(\eps^2),\qquad
 \phi_\eps=\phi_0+\eps\eta+\tfrac12\eps^2w+o(\eps^2).
 \label{eq:factorial}
\end{equation}
An angular mean uses probability measure on $S^3$. The first angular
factor is $Y=n_1^2-n_2^2$, so
$\av{Y}=0$, $\av{Y^2}=1/6$, and $\Delta_\kappa Y=-8Y$.

The following are explicit analytic premises, not theorems reproved by
the interval calculation: the regular background and its stated collar
bounds; the invariant dotted-space weighted constraint inverse and its
nonlinear implicit-function construction; a common fixed-shell development
with the retained subsidiary and boundary uniqueness statements; its $C^2$
parameter dependence with the lower spacetime regularity needed for traces
and observables; and identification of the actual first master and harmonic
gauge fields with that family's first derivative. The companion technical
records specify these premises and the fixed first seed \cite{companion}.
The common solution slab contains the target with a strict time margin.
For each sufficiently small fixed family, the anchored intrinsic normal
foliation exists on the needed compact interval after qualitative parameter
shrinkage. No numerical no-caustic or parameter-radius estimate is supplied.

\begin{proposition}[Conditional response at a fixed proper clock]
Under these premises, choose the explicit mean preparation of
Section~\ref{sec:preparation}, keep the retained first packet, and define
the anchored observables of Section~\ref{sec:proper}. At
$\tau_*=(5/4)\rho_b\ell$, their second derivatives obey
\begin{align}
 M_2(\tau_*)&\in[1.1296,1.9779]\times10^{30},\nonumber\\
 v_2(\tau_*):=V_2(\tau_*)/V_0(\tau_*)&\in[-5.4520,-2.6521]\times10^{16},
 \label{eq:main}\\
 \Theta_2(\tau_*)&\in[-9.0208,-4.4458]\times10^{21}.\nonumber
\end{align}
The intervals are outward decimal displays of exact rational enclosures.
Every fixed admissible choice of the remaining radial weights has this
same second mean response. This does not identify the full finite-parameter
families or give a uniform admissible parameter interval over those choices.
\end{proposition}

\section{An explicit compact second mean}\label{sec:preparation}

Write $\rho_0$ for the background spatial volume density, distinct from
$\rho$. With the action normalization $R/2-|\dd\phi|^2/2-\mathcal V$,
the initial first canonical data in a background orthonormal spatial frame
are
\begin{equation}
 \gamma_1=\operatorname{diag}(-4u,2u,2u,2u)Y,\quad
 \pi_1/\rho_0=\operatorname{diag}(-3s/2,0,0,0)Y,\quad
 \phi_1=vY,\quad \varpi_1/\rho_0=jY, \label{eq:first}
\end{equation}
where $s=u_t/\rho$, $j=v_t/\rho$. The actual initial packet is
\begin{align}
 Q(z)&=\exp\!\left(4-\frac{\ell^2}{(z-\ell)(2\ell-z)}\right)
 &&(\ell<z<2\ell), \label{eq:bump}\\
 \alpha&=(\rho/\rho_b)^{3/2}p/p_b,\quad
 \kappa_r=\alpha_z/\alpha=\rho(5a/2-\mathcal V'/p),\nonumber\\
 u&=\alpha Q/\rho^2,\quad u_t=e^{-z}u,\quad
 v=\frac{3\alpha(Q_z+\kappa_rQ)}{\rho^3p},\quad
 v_t=\frac{3\alpha e^{-z}[Q_z+(\kappa_r-1)Q]}{\rho^3p}.\nonumber
\end{align}
Set $Q=0$ elsewhere; its endpoints are flat. In particular
$pv=-3(u_y+2au)$ and $pj=-3(s_y+3as)$.

Let $Q_H$ and $Q_M$ be the monopole Hessians of the canonical Hamiltonian
constraint and the radial momentum constraint divided by $\rho_0$.
Write the latter as $Q_M^{\rm mom}$. Their explicit expressions are
\begin{align}
 Q_H={}&-2u_y^2-24auu_y-2uu_{yy}
 -(16a_y+32a^2+8p^2/3)u^2-4u^2/\rho^2-s^2\nonumber\\
 &-(j^2+v_y^2+\mathcal V''v^2)/6
 -(4/3)puv_y-(4/3)v^2/\rho^2, \label{eq:sources}\\
 Q_M^{\rm mom}={}&-2s u_y+j(v_y+4pu)/3.\nonumber
\end{align}
These are Hessians, not coefficients already divided by two. If $P$ is
the linearized constraint operator, the required second equation is
$P x_2=-D^2\widehat{\mathcal C}[v_{\rm seed},v_{\rm seed}]$.

With $y_L=y(2\ell)<y_U=y(\ell)$, define
\begin{align}
 I&=\int_{y_L}^{y_U}\rho^4 Q_H\,\dd y,&
 J&=\int_{y_L}^{y_U}\rho^2 Q\,\dd y>0,& f&=\frac{2I}{9J}Q,
 \label{eq:IJ}\\
 C[f]&=-\tfrac32(f_{yy}+3af_y)-\mathcal V f,&
 \xi(y)&=-\frac{\displaystyle\int_{y_L}^{y}\rho^4(Q_H+C[f])\,\dd s}
                    {\rho^4p^2}.&&\label{eq:xi}
\end{align}
The mean canonical datum $B_{\rm mean}$ is
\begin{align}
 h&=f\gamma_0+\Lie_{\xi\partial_y}\gamma_0,\qquad
       \phi_2=\varpi_2=0,\nonumber\\
 \tau:=\pi_2/\rho_0&=\operatorname{diag}
       (0,-Q_M^{\rm mom}/(6a),-Q_M^{\rm mom}/(6a),-Q_M^{\rm mom}/(6a)).
       \label{eq:B}
\end{align}
The displayed metric components are $h_{yy}=f+2\xi_y$ and
$h_{\rm angular}=f+2a\xi$.

The radial identities proving the construction are
\begin{align}
 P_g(f\gamma_0)&=C[f],&
 P_g(\Lie_{\xi\partial_y}\gamma_0)&=\rho^{-4}(\rho^4p^2\xi)_y,
 \nonumber\\
 \int\rho^4 C[f]\,\dd y&=-\tfrac92\int\rho^2f\,\dd y,&
 P_m(C,D)&=-2[C_y+3a(C-D)]. \label{eq:rightinverse}
\end{align}
Thus the primitive in \eqref{eq:xi} has zero total integral and vanishes
on both exterior sides. Flatness gives a smooth compact extension, and
\eqref{eq:B} solves $P B_{\rm mean}=-P_0D^2\widehat{\mathcal C}[v,v]$.
The denominator bounds $\rho\ge129$, $|p|\ge3/25000$, and
$a\ge111/1000$ hold on the support. The sign and factor
$-Q_M^{\rm mom}/(6a)$ are essential.

For completeness, let $\mathcal A$ be the retained radial weighted
correction map and $\omega$ its constraint weight. Apply the same local
implicit-function theorem to
\begin{equation}
 x_\eps=x_0+\eps v_{\rm seed}+\tfrac12\eps^2B_{\rm mean}
             +\mathcal A U(\eps),\qquad U(0)=U'(0)=0. \label{eq:ift}
\end{equation}
The derivative in $U$ is the original isomorphism
$L=\omega^{-2}P\mathcal A$. Smooth compact polynomial trial terms only
alter local constants. Haar averaging $P_0$ on the full tensor spaces
commutes with the radial operators and weights, is bounded on the dotted
closures, and maps the reflection-invariant domain into itself. Hence
$LP_0U''(0)=0$, so injectivity gives $P_0U''(0)=0$. Individual rotations
need not preserve that reflection-invariant domain. Consequently the exact
curve has the prescribed second mean for every remaining admissible radial
weight choice. Other angular second data and higher parameter derivatives
are not thereby fixed.

\section{A conditioned second harmonic jet}

The canonical-to-physical map is
\begin{equation}
 K_{ij}=\frac{2}{\rho_\gamma}
 \left(\gamma_{ia}\gamma_{jb}\pi^{ab}
       -\tfrac13\gamma_{ij}\gamma_{ab}\pi^{ab}\right),\qquad
 \Pi=\varpi/\rho_\gamma,\qquad K=\tfrac12\Lie_n\gamma.
 \label{eq:canonical}
\end{equation}
At the time-symmetric background it gives
$\av{K_{2,yy}}=Q_M^{\rm mom}/(3a)+6us$,
$\av{K_{2,\rm angular}}=-us$, and $\av{\Pi_2}=-uj/3$.
Choose initial lapse $N_\eps=\rho(1+\eps uY)$ and zero shift.
For $R^A=g^{BC}\widehat\Gamma^A_{BC}$ with
$\widehat\Gamma=\Gamma[g_0]$, the exact initial relative-harmonic
completion is
\begin{equation}
 \dot N=N^2\operatorname{tr}_\gamma K-N^3R^t,\qquad
 \dot\beta^I=N^2(\Gamma_\gamma^I-R^I)-N\gamma^{IJ}\partial_JN.
 \label{eq:gaugeinitial}
\end{equation}
The perturbed inverse metric contracts the reference connection. Together
with $\dot\gamma=2NK$ and $\dot\phi=N\Pi$, this fixes all gauge
velocities. The physical constraints then supply the vanishing initial
normal subsidiary derivative through the reduced equations; those
accelerations are not extra freely chosen Cauchy data.

The spatial term $2\xi_y$ can have size about $10^{39}$ in this
preparation. Evolving the original angular trace requires cancellation
against a large radial contribution in the shell boundary lift. We use
a pure second generator $Z$, initially
\begin{equation}
 (Z^y,Z^t,Z^y_t,Z^t_t)=(-\xi,0,0,a\xi),\qquad Z^{\rm angular}=0,
 \label{eq:Zdata}
\end{equation}
and solve $(\Box_0+\operatorname{Ric}_0)Z=0$, with
$Z^y=0$ and $\partial_yZ^t=0$ at the shell. These are the monopole
normal Dirichlet and tangential Neumann rows. They are derived directly
for the monopole vector; an angular-gradient amplitude is absent.
The smooth compatible two-component problem has a standard reflected
Volterra construction. The companion check bounds its contraction
constant on the retained triangle by $3/250$.

The second transformation is
\begin{equation}
 \bar k=k+\Lie_Zg_0,\qquad \bar w=w+Z\phi_0. \label{eq:secondgauge}
\end{equation}
It removes the spatial Lie term and gives $\bar\phi_2=-p\xi$.
It leaves second canonical momenta unchanged because their background
values vanish; naturality of \eqref{eq:canonical} preserves the physical
conversion. The chosen time derivative in \eqref{eq:Zdata} keeps the
second lapse and shift values fixed. Omitting the new scalar displacement
would change the physical data and violate the linearized constraint
cancellation.

Use the background frame $e_y=\partial_y$, $e_0=\rho^{-1}\partial_t$,
$e_a=(\rho\cosh t)^{-1}E_a$, and mean order
$\mathbf U=(K_0,L_0,N_0,J_0,w_0)$ for the $yy,y0,00$, angular-isotropic,
and scalar slots. The positive angular norm has Gram matrix
$\operatorname{diag}(1,2,1,3,1)$. In this conditioned representative,
\begin{align}
 \mathbf U_0&=(f,0,-u^2/3,f,-p\xi), \label{eq:U0}\\
 \mathbf V_0:=\partial_x\mathbf U|_0
 &=\ell\begin{pmatrix}
 2\rho Q_M^{\rm mom}/(3a)+32uu_t/3\\
 \rho[-f_y+af-13uu_y/3-au^2/3]\\
 -2\rho Q_M^{\rm mom}/(3a)+4uu_t/3\\
 -4uu_t/3\\0
 \end{pmatrix},\label{eq:V0}\\
 \partial_r\mathbf U_0&=-\ell\rho(f_y,0,-2uu_y/3,f_y,-p_y\xi-p\xi_y).
 \label{eq:U0r}
\end{align}
In particular $N_0$ is a covariant timelike frame component, with no
additional sign flip. The $g_{ty}$ velocity includes
$2(\gamma_1)_{yy}\dot\beta_1^y$; the $g_{tt}$ value includes the
quadratic lapse square. The vanishing scalar time velocity is the exact
cancellation $\rho\Pi_2+2N_1\Pi_1=0$.

This is a change of second jet of the same prepared physical family.
The flow of $Z$ at parameter $\eps^2/2$ preserves the shell and its
initial anchor. It need not be a nonlinear harmonic transformation for
finite $\eps$. Only $dC[\Lie_Zg_0]=0$ is required for the second jet.
The first field and all its quadratic sources remain unchanged. The
retained parameter regularity permits the needed second-jet and trace
operations; no new lossless Sobolev composition theorem is asserted.

\section{The complete five-mean estimate}

Rotational equivariance of the linearized background bulk and shell
operators makes all five invariant slots a closed system, even for second
data with additional angular components. Projecting the exact family's
second equations and applying the retained mixed linear uniqueness
identifies their mean with the solution from
\eqref{eq:U0}--\eqref{eq:U0r}. The operator, physical shell rows, and
normal gauge row are retained together:
\begin{equation}
 \mathbf U_{xx}-\mathbf U_{rr}
 =A_0\mathbf U_x+B_0\mathbf U_r+D_0\mathbf U+F_0,
 \quad L_0|_{r=0}=0,\quad P_N\mathbf U_r=M_0P_N\mathbf U+G_0.
 \label{eq:meanpde}
\end{equation}
Here $P_N$ selects the other four slots. The actual forcing includes
the reduced Einstein/scalar Hessians and all physical and gauge shell
sources, not a selected-sign replacement.

Set $T=5/4$, $\Delta=T-1=1/4$, and
\begin{equation}
 D_T=\{0\le x\le T,\ 0\le r\le T-x\},\qquad
 \norm{\mathbf U}_X=\max\left\{sup_{D_T}|\mathbf U|_H/3,
      \sup_{D_T}|\mathbf U_x|_H,\sup_{D_T}|\mathbf U_r|_H\right\}.
 \label{eq:norm}
\end{equation}
The mean initial data and first quadratic sources vanish in the flat
wedge $x+r\le1$. Thus the linear mean response vanishes there. This
does not impose that same wedge on every finite-parameter nonlinear field.
For compatible differences, integration from the wedge gives
$|\mathbf U|_H\le\Delta\norm{\mathbf U}_X$ and shell trace
$|\mathbf U(x,0)|_H\le(x-1)_+\norm{\mathbf U}_X$.

Write $\mathcal K$ for the reflected bulk Duhamel operator and
$\mathcal L$ for the Robin lift. Their bounds are
$\norm{\mathcal KF}_X\le T\sup|F|_H$ and
$\norm{\mathcal LG}_X\le\sup|G|_H$.
The retained coefficient caps
$\norm{A_0},\norm{B_0}\le10^{-4}$,
$\norm{D_0}\le10^{-6}$ and $\norm{M_0}\le1/2000$ yield
\begin{equation}
 q=T(2\cdot10^{-4}+\Delta10^{-6})+\Delta/2000
   =\frac{1201}{3200000}=0.0003753125<1. \label{eq:q}
\end{equation}
If $\mathbf B$ is the reflected free wave for the \emph{new} initial
pair, the exact total mean satisfies
$\mathbf U=\mathbf B+\mathcal T\mathbf U+\mathbf S$.
The entire forcing history is charged once using
\begin{equation}
 F_H=17387513156573,\quad G_H=3851400,\qquad
 S_X=TF_H+G_H=\frac{86937581188465}{4}. \label{eq:sourcecap}
\end{equation}
For new initial norm caps $D_0^{\rm init},D_r^{\rm init},D_v^{\rm init}$,
\begin{equation}
 B_X\le\max\{(D_0^{\rm init}+TD_v^{\rm init})/3,
                    D_r^{\rm init}+D_v^{\rm init}\},\qquad
 \norm{\mathbf U-\mathbf B}_X\le E_X:=\frac{qB_X+S_X}{1-q}.
 \label{eq:resolvent}
\end{equation}
This bounds all feedback iterations and all internal mean couplings.
No old total cap, initial accuracy allocation, or truncated-iteration tail
is substituted. Whole-cell interval evaluation with flat-tail treatment
gives the outward displays
\begin{center}
\begin{tabular}{@{}lc@{}}\toprule
Quantity & Certified upper display\\\midrule
$D_0^{\rm init}$ & $1.9944\times10^{30}$\\
$D_r^{\rm init}$ & $1.7478\times10^{31}$\\
$D_v^{\rm init}$ & $2.7533\times10^{22}$\\
$B_X$ & $1.7478\times10^{31}$\\
$E_X$ & $6.5620\times10^{27}$\\\bottomrule
\end{tabular}
\end{center}
The large coefficients reflect the chosen parameter normalization; they
are not finite-amplitude field values. The construction bounds a
continuous response rather than numerically time-stepping a full field.

\section{Proper scalar mean and intrinsic volume}\label{sec:proper}

Anchor the shell's intrinsic future normal-geodesic foliation at its
registered initial spatial section. Let $h_\eps(\tau)$ be the induced
spatial metric in that foliation and define
\begin{equation}
 V_\eps(\tau)=\int_{S^3}\dd\mu_{h_\eps},\qquad
 \Theta_\eps=\partial_\tau\log V_\eps,\qquad
 M_\eps=V_\eps^{-1}\int_{S^3}\phi_\eps\,\dd\mu_{h_\eps}.
 \label{eq:observables}
\end{equation}
The observers are intrinsic shell geodesics, not an assumed material flow
or ambient geodesic congruence. The background has
$A(\tau)=\rho_b\cosh(\tau/\rho_b)$,
$H=A_\tau/A$, and $V_0=2\pi^2A^3$.
All first integrated coefficients vanish by the quadrupole mean.

To exhibit the required first corrections, write the first shell frame
metric as $m_{00}=-2XY$, $m_{0a}=V D_aY$, and
$m_{ab}=2(ZY\kappa_{ab}+dD_aD_bY)$, and write the first scalar as
$\psi Y$. Here $V$ is a first metric amplitude only in these formulas,
not the shell volume. Set
$C_\tau=X$, $C(0)=0$, and $E=Z-HC$.
The normalized proper mean is
\begin{equation}
 M_2=w_0-C\psi_\tau/3+(E-8d/3)\psi. \label{eq:M2}
\end{equation}
The last term is the perturbed measure correction. The $d$ contribution
is retained. Angular transport cancels only after scalar pullback and
measure are combined, not by dropping both indiscriminately.

Let $Q_{\rm corr}$ bound the last two terms of \eqref{eq:M2}. The actual
first shell two-jet is enclosed on the whole interval $1\le x\le T$,
giving $Q_{\rm corr}<3.7987$. Even reflection and zero initial scalar
velocity give $B_w(x,0)=w_{\rm init}(x)=-p\xi|_{r=x}$. Hence
\begin{equation}
 |M_2(x)-w_{\rm init}(x)|\le(x-1)E_X+Q_{\rm corr}\quad(1\le x\le T).
 \label{eq:Merror}
\end{equation}
Its linearly weighted and unweighted history error bounds are, respectively,
$\Delta^3 E_X/6+\Delta^2Q_{\rm corr}/2$ and
$\Delta^2 E_X/2+\Delta Q_{\rm corr}$. Exact positive cell weights integrate
the continuous initial scalar boxes. No assumption of pointwise positivity
at every arbitrarily early contact time enters these history bounds.

The exact shell Gauss contraction of \eqref{eq:model}--\eqref{eq:shell}
is
\begin{equation}
 R(\gamma)=(D\phi)^2+F(\phi),\qquad
 F=2\mathcal V-\tfrac14(\sigma')^2+\tfrac13\sigma^2. \label{eq:gauss}
\end{equation}
This standard geometric identity removes the electric Weyl trace; it
does not determine every induced tensor equation. Define
$v_2=V_2/V_0$. The volume-weighted second identity in the anchored
proper-time foliation is
\begin{align}
 (v_2)_{\tau\tau}+4H(v_2)_\tau-2v_2/A^2
    &=\tfrac12F'M_2+Q_1,\qquad v_2(0)=(v_2)_\tau(0)=0,
 \label{eq:volumeODE}\\
 Q_1={}&F''\psi^2/12-\psi_\tau^2/6+4\psi^2/(3A^2)-5E^2/(3A^2)
 \nonumber\\
 &+E_\tau^2-16E_\tau D_\tau/3+8D_\tau^2/3. \label{eq:Q1}
\end{align}
Here $D$ is the synchronous angular Hessian coefficient, with
$D_\tau=d_\tau-V/A-C/A^2$. Every term is retained and bounded in absolute
value; the resulting cap is $|Q_1|<7.8317\times10^{19}$.
In \eqref{eq:volumeODE} the scalar source is the full proper mean,
not the coordinate coefficient $w_0$.

One way to check the curvature coefficient is to vary the functional
$\mathcal J(h)=\int R_3\dd\mu-(2/A^2)\int\dd\mu$ at the round sphere.
It is critical. For the coordinate metric variation
$h_1=2A^2(EY\kappa+D\operatorname{Hess}Y)$, its normalized Hessian is
$V_0^{-1}D^2\mathcal J[h_1,h_1]=10E^2/(3A^2)$; a pure spatial
diffeomorphism contributes zero.
The weighted second spatial curvature is therefore
$-4v_2/A^2+10E^2/(3A^2)$. The temporal part contributes
$2(v_2)_{\tau\tau}+8H(v_2)_\tau
 +2\av{\operatorname{tr}K_1^2-\theta_1^2}$, with
\begin{equation}
 \av{\operatorname{tr}K_1^2-\theta_1^2}
 =-E_\tau^2+16E_\tau D_\tau/3-8D_\tau^2/3.
\end{equation}
Combining this with
$\av{(D\phi)^2}_2=-\psi_\tau^2/3+8\psi^2/(3A^2)$ proves
\eqref{eq:volumeODE}--\eqref{eq:Q1}, including the derivative convention.

In scaled time the exact Green kernel is
\begin{equation}
 k(T,x)=\frac{\sinh(2\ell(T-x))}{2\ell}
             \frac{\cosh^2(\ell x)}{\cosh^2(\ell T)}. \label{eq:Green}
\end{equation}
It follows from the homogeneous solutions
$\tanh(\ell x)$ and $1+\tanh^2(\ell x)$. It is positive, with
$0\le k(T,x)\le T-x$. On the interval used here,
$|k/(T-x)-1|<10^{-11}$ and $|k_T-1|<10^{-11}$, interpreted continuously
at the endpoint. These uniform errors are charged to the history bounds.
Consequently
\begin{align}
 v_2(T)&=(\rho_b\ell)^2\int_1^T
       k(T,x)[F'M_2(x)/2+Q_1(x)]\,\dd x, \label{eq:integralvolume}\\
 \Theta_2(T)=(v_2)_\tau(T)&=\rho_b\ell\int_1^T
       k_T(T,x)[F'M_2(x)/2+Q_1(x)]\,\dd x. \label{eq:integraltheta}
\end{align}
Here $F'$ is evaluated at the constant background shell scalar. Its
retained interval is negative. Equations \eqref{eq:Merror} and
\eqref{eq:integralvolume}--\eqref{eq:integraltheta}, with the full
source budgets, give \eqref{eq:main}. In particular $\Theta_2$ is the
first proper-time derivative of $v_2$, not an acceleration.

\section{Interpretation, verification, and relation to prior work}

At the chosen positive time the background expansion is
\begin{equation}
 \Theta_0(\tau_*)=\frac{3}{\rho_b}\tanh(\ell T)>0.
\end{equation}
Because the first integrated coefficients vanish, $C^2$ parameter
dependence gives, at this same anchored proper clock,
\begin{equation}
 \frac{V_\eps}{V_0}=1+\tfrac12\eps^2v_2+o(\eps^2),\qquad
 \Theta_\eps=\Theta_0+\tfrac12\eps^2\Theta_2+o(\eps^2).
\end{equation}
The strict signs in \eqref{eq:main} therefore imply, for each fixed
admissible family and sufficiently small nonzero $\eps$ of either sign,
\begin{equation}
       0<V_\eps(\tau_*)/V_0(\tau_*)<1,\qquad
       0<\Theta_\eps(\tau_*)<\Theta_0(\tau_*). \label{eq:qualitative}
\end{equation}
No particular nonzero $\eps$ is certified. The allowed interval depends
on the family and its remainder control, and need not be uniform in the
remaining weights. The background is still expanding in
\eqref{eq:qualitative}; a negative second change is not total contraction.

The numerical statements use whole-cell outward interval arithmetic and
exact rational endpoints, including the packet's flat tails. The initial
campaign uses the registered $N=1024$ base grid, with 994 recorded cells:
992 interior cells and two grouped flat-tail regions. The conditioned adapter is
independently checked against every supplied cell, and the complete source
history enters once through \eqref{eq:sourcecap}. Normal and optimized
runs agree exactly. A separate reconstruction of the response budget
matches all three final rational intervals without importing the response
producer. These checks authenticate and audit the stated finite
calculation; they do not constitute proofs of the retained analytic
existence assumptions or external peer review. The accompanying materials
provide the proofs, arithmetic adapters, source pins, raw boxes, and
canonical reports \cite{companion}.

The calculation uses established geometric and analytic machinery.
Bruni et al.\ \cite{Bruni} give the second-order gauge law, including
its factorial convention and first-generator cross terms; Mars
\cite{Mars} treats perturbations of hypersurfaces and the distinction
between spacetime gauge and hypersurface identification. Transport of the
shell and its anchor is necessary here. The Gauss--Codazzi projection
used in brane geometry \cite{SMS} is standard; constant-tension or
constant-bulk matter reductions cannot be imported unchanged into
\eqref{eq:gauss}. Buchert's averaging identities \cite{BuchertI,BuchertII}
establish the expansion-variance/shear combination and the averaging
commutator. Their dust or fluid evolution assumptions are not imposed on
our intrinsically geodesic shell observers. The reflected-kernel
contraction and Green-function calculation are standard linear methods.

The contribution is the model-specific connection between an explicit
compact second mean, its conditioned five-field response, and these
anchored observable enclosures. This is a focused comparison with those
primary frameworks, not an exhaustive novelty survey or a new general
implicit-function theorem, averaging law, or Einstein equation. The
uncomputed higher angular sectors are not needed for this invariant
projection, but no complete angular waveform is supplied. The large
second coefficients have no SI calibration here. No quantitative
finite-amplitude continuation, hot-matter or radiation history, physical
packet-generation mechanism, or Big-Bang-origin conclusion follows.

\paragraph{Acknowledgment and review status.}
The manuscript and accompanying calculations were prepared with AI
assistance, including internal algebraic and reproducibility review.
This working draft is supplied for author and specialist review. Internal
checks are not external peer review.

\begin{thebibliography}{9}
\footnotesize
\setlength{\itemsep}{3pt}
\bibitem{Bruni} M. Bruni, S. Matarrese, S. Mollerach, and S. Sonego,
\emph{Perturbations of spacetime: gauge transformations and gauge invariance
at second order and beyond}, Classical and Quantum Gravity \textbf{14}
(1997), \href{https://arxiv.org/abs/gr-qc/9609040}{arXiv:gr-qc/9609040}.
\bibitem{Mars} M. Mars, \emph{First and Second Order Perturbations of
Hypersurfaces}, Classical and Quantum Gravity \textbf{22} (2005),
\href{https://arxiv.org/abs/gr-qc/0507005}{arXiv:gr-qc/0507005}.
\bibitem{SMS} T. Shiromizu, K. Maeda, and M. Sasaki,
\emph{The Einstein Equations on the 3-Brane World},
\href{https://arxiv.org/abs/gr-qc/9910076}{arXiv:gr-qc/9910076}.
\bibitem{BuchertI} T. Buchert, \emph{On Average Properties of Inhomogeneous
Fluids in General Relativity I: Dust Cosmologies},
\href{https://arxiv.org/abs/gr-qc/9906015}{arXiv:gr-qc/9906015}.
\bibitem{BuchertII} T. Buchert, \emph{On Average Properties of Inhomogeneous
Fluids in General Relativity II: Perfect Fluid Cosmologies},
\href{https://arxiv.org/abs/gr-qc/0102049}{arXiv:gr-qc/0102049}.
\bibitem{companion} R. Maldonado, \emph{HDBLAST conditional technical
records and reproducibility materials} (2026),
\href{https://github.com/maldonado-research/HDblast}{repository}.
Retained records: exact localized initial data (14 October), common-gauge
response (16 October), exact-family transfer (31 October), and intrinsic
shell budget (1 November). The associated prepared-response materials
contain the explicit preparation, both initial-gauge adapters, and the
exact final scalar/volume/expansion report used here.
\end{thebibliography}
\end{document}
